Participants have to identify as many stars as possible in the field of a
celestial photograph using the telescope and CCD provided by the committee.
Participants have to choose one out of the five recommended regions in the
sky listed below. Then, point the telescope to the direction of the selected
sky region. Take three photographs with different exposure times and record
the images of the sky by the CCD camera. Save the observational data.
Transfer the data to the printing facilities to print out the result. Ask
the technical assistant for a help. Choose the best prints-out and use the
image to identify the stars in the field of observations. The following
procedures are,

\begin{enumerate}
\item Choose only one out of the following five recommended regions to the
  directions of (marked by the following bright clusters):
  \begin{description}
  \item[i.] M7 ($\alpha = 17^{\mathrm{h}}53.^{\mathrm{m}}3$,
    $\delta = -34^\circ 46.'2$)
  \item[ii.] M8 ($\alpha = 18^{\mathrm{h}}04.^{\mathrm{m}}2$,
    $\delta = -24^\circ 22.'0$)
  \item[iii.] M20 ($\alpha = 18^{\mathrm{h}}02.^{\mathrm{m}}4$,
    $\delta = -22^\circ 58.'7$)
  \item[iv.] M21 ($\alpha = 18^{\mathrm{h}}04.^{\mathrm{m}}2$,
    $\delta = -22^\circ 29.'5$)
  \item[v.] M23 ($\alpha = 17^{\mathrm{h}}57.^{\mathrm{m}}0$,
    $\delta = -18^\circ 59.'4$)
  \end{description}
  You may not change your choice.
\item Type in your country code, your student code, and your choice of
  region into the answer page provided in the computer.
\item Point the telescope to the chosen cluster by using telescope
  controller. If necessary you may move the telescope slightly to get the
  best position in the frame of the CCD by checking the display of CCDops
  software.
\item Display the region in The Sky map software provided in the computer,
  to confirm that the telescope is pointed to the selected object in the
  sky. You may change field of view of the sky chart.
\item You may invert the background images into white color, as in the chart
  mode. Copy and paste the sky chart from The Sky into the answer page. Use
  ``Ctrl c'' and ``Ctrl v'' buttons from the keyboard, respectively.
\item Type the equatorial coordinates of the centre of that object in the
  answer page as indicated in The Sky.
\item Take three photographs of the chosen object by using the attached CCD
  camera and CCDops software, with various exposure times. Choose exposure
  time in the range between 1 and 120 seconds. Image is automatically
  subtracted with dark frame with the same exposure time.
\item You must invert the background images into white color. Copy and paste
  images from CCDops into the answer page. Use ``Ctrl c'' and ``Ctrl v''
  buttons from the keyboard, respectively.
\item Save your answer page into hard disk in a Microsoft Word file format.
\item Print your answer page which consists of the photographs and the
  corresponding sky chart.
\item Go to the identification room and bring with you the prints-out of the
  sky chart and the photographs. Ask some help from technical assistant, if
  necessary.
\item Choose the best out of the three printed images and identify as many
  objects as possible on it.
\item Use the assigned computer and The Sky software to identify objects.
  Type your identification to your answer page.
\item Type on the answer page the names (or catalogue number), RA, Dec, and
  magnitude of each identified star and put the sequential number on the
  photograph. Make sure that you list the stars on the answer page following
  the same order and number as on the photograph.
\item Estimate the limiting magnitude of the photograph you choose,
  empirically.
\end{enumerate}
